\section{Extending a connection through a fresh region}
\label{sec:extension}

The construction adapts the shell, separated-seed, and conditional-target
strategy of Kozma and Nitzan~\cite[Section~4, Lemma~10]{KozmaNitzan2024}.
We use the hyperedge failure comparison for both models and define relay
quality using the interior bits alone.

We retain the parameters chosen in Lemma~\ref{prelim:scales}.
A rectangle is a nonempty set of the form
\[
 B=\prod_{h=1}^3([a_h,b_h]\cap\mathbb Z),
 \qquad a_h,b_h\in\mathbb Z,
\]
and $B^{(j)}$ denotes its expansion by $j$ in every coordinate direction.
Bits internal to a vertex set mean sites in that set, or bonds with both
endpoints in it, according to the model.

\begin{lemma}[Extension estimate]
\label{ext:estimate}
Let $G$ be a finite subgraph of $\mathbb Z^3$ containing a rectangle $D$
and every lattice edge internal to $D$.
Give its bits independent laws, with parameter $q$ on all bits internal
to $D$ and arbitrary parameters, including $0$ and $1$, elsewhere.
Let $o\in V(G)$ and $T\subseteq V(G)$.
Suppose that a rectangle $B$ satisfies
\[
 B^{(R)}\subseteq D,
 \qquad o\notin B^{(R)},
\]
and that for every $u\in B^{(R)}$ there is a radius
$\ell\in\{2r,2r+1,10r\}$ such that $\Lambda_\ell(u)\subseteq D$
and some quarter-face of this cube is contained in $T$.
Then, for connections in $G$,
\begin{equation}
 \mathbb P(o\leftrightarrow B,\ o\not\leftrightarrow T)<\eta.
 \label{ext:conclusion}
\end{equation}
\end{lemma}

\begin{proof}
We first choose a deterministic layer at which a connection to $B$
usually supplies many possible entry points. We then try disjoint open
seed boxes at those points, and use Corollary~\ref{cmp:failure} to pass
from a successful seed to $T$.

\emph{Choosing a layer.}
The disjoint closure events used here also appear in Grimmett and
Marstrand~\cite[Section~3, Lemma~3]{GrimmettMarstrand1990}.
For $m\le j<R$, expose only the bits internal to the induced graph
outside $B^{(j)}$. Let $Y_j$ be the exterior neighbors of $B^{(j)}$
reachable from $o$ by an open path in this outside graph, and set
$N_j=|Y_j|$.
Every exterior neighbor $y$ of a rectangle has a unique neighbor
$z_y$ inside it. Both endpoints and their edge lie in
$B^{(j+1)}\subseteq D$.
Since $o$ is outside $B^{(R)}$, any path from $o$ to $B$ has a first
entry into $B^{(j)}$, and hence $N_j\ge1$.

On $1\le N_j<L$, consider the event that all bonds $\{y,z_y\}$,
$y\in Y_j$, are closed in the bond model, or that all corresponding
sites $z_y$ are closed in the site model.
Given the outside bits, these are unexposed parameter-$q$ bits.
There are at most $N_j$ distinct bits; repeated sites $z_y$ can only
increase the closure probability. Thus its conditional probability
is at least
\[
 (1-q)^{N_j}\ge (1-p)^L=\alpha.
\]
Such a closure prevents every first entry into $B^{(j)}$.
In particular, $N_i=0$ for all $i<j$, because every exterior neighbor
of $B^{(i)}$ belongs to $B^{(j)}$.
The events consisting of $1\le N_j<L$ and this closure are therefore
pairwise disjoint as $j$ varies. It follows that
\[
 \sum_{j=m}^{R-1}\mathbb P(o\leftrightarrow B,\ N_j<L)
 \le \sum_{j=m}^{R-1}\mathbb P(1\le N_j<L)
 \le \alpha^{-1}.
\]
Fix a deterministic $j\in\{m,\ldots,R-1\}$ for which
\begin{equation}
 \mathbb P(o\leftrightarrow B,\ N_j<L)
 \le \frac{1}{\alpha(R-m)}.
 \label{ext:layer}
\end{equation}
Write $C=B^{(j)}$. Let $\xi$ be the configuration of bits internal
to $C$, and let $\zeta$ be the configuration of bits internal to the
induced graph outside $C$. Thus $Y_j$ is determined by $\zeta$.
For bonds, the remaining bits are the bonds crossing between $C$
and its complement. The configurations $\xi$, $\zeta$, and these
crossing bits are independent. All bits of $\xi$ and all crossing
bits have parameter $q$, since they lie in $B^{(j+1)}\subseteq D$.
For sites, $\xi$ and $\zeta$ already partition the bits.

We will use these two configurations for different purposes.
The exterior $\zeta$ selects entrances that $o$ has reached, whereas
$\xi$ determines which relay centers have high conditional probability
of reaching $T$. We first estimate relay quality for each fixed entrance,
then use independence to apply that estimate to the entrances selected
by $\zeta$.

\emph{Defining good relays from $\xi$.}
For every exterior neighbor $y$ of $C$, choose $v_y$ so that
\[
 U_y=\Lambda_m(v_y)\subseteq C,
 \qquad z_y\in U_y,
 \qquad \|v_y-y\|_\infty\le m+1.
\]
Explicitly, clamp each coordinate of $z_y$ to the corresponding
coordinate interval of $C$ shrunk by $m$ at both ends. Those shrunken
intervals are nonempty because $j\ge m$.
Let $O_y$ be the event that every internal bit of $U_y$ is open.
Then $O_y$ is determined by $\xi$ and $\mathbb P(O_y)=q^M$.
Define
\[
 g_y=\mathbb P(v_y\leftrightarrow T\mid\xi),
 \qquad
 A_\xi=\{v_y:y\text{ is an exterior neighbor of }C,
                         \ g_y>1-\tau\}.
\]
In this definition the exterior configuration and the crossing bits
are still averaged over according to their original product law.
In particular, $A_\xi$ is determined by $\xi$ alone.

For a fixed $y$, apply the geometric hypothesis at $u=v_y$.
The quarter-face estimate \eqref{prelim:quarterface} supplies an
increasing event of probability greater than $1-\gamma$ connecting
$U_y$ to $T$ inside the specified cube. On its intersection with $O_y$,
the center $v_y$ also connects to $T$. Positive association gives
\[
 \mathbb P(O_y,\ v_y\not\leftrightarrow T)\le q^M\gamma.
\]
Since $O_y$ is determined by $\xi$, it follows that
\begin{equation}
 \begin{split}
 \tau\,\mathbb P(O_y,\ g_y\le1-\tau)
 &\le \mathbb E\big[\mathbf 1_{O_y}(1-g_y)\big]\\
 &=\mathbb P(O_y,\ v_y\not\leftrightarrow T)
 \le q^M\gamma.
 \end{split}
 \label{ext:relay-quality}
\end{equation}
This bounds the chance that a fixed seed is open but its center fails
the defining condition for $A_\xi$.

\emph{Selecting reachable entrances using $\zeta$.}
Condition now on $\zeta$. Whenever $N_j\ge L$, choose deterministically
$K$ members of $Y_j$ at pairwise sup distances greater than $4m+2$.
This is possible by greedy selection: each chosen vertex excludes at
most $(8m+5)^3$ candidates, while $L=K(8m+5)^3$.
The corresponding centers have pairwise distances greater than $2m$,
so their boxes $U_y$ are vertex-disjoint.

Call a selected entrance successful if $O_y$ holds and, for bonds,
the crossing bond $\{y,z_y\}$ is open. For sites, $y$ is already
open because it is reachable in the outside graph, and the open
seed contains $z_y$. In either model a success therefore gives an
actual open path from $o$ to $v_y$.

Given $\zeta$, these successes are independent Bernoulli trials with
common parameter
\[
 c=\begin{cases}
 q^{M+1},&\text{for bonds},\\
 q^M,&\text{for sites}.
 \end{cases}
 \qquad c\ge\rho_0.
\]
Indeed, the boxes use disjoint sets of bits of $\xi$, and the selected
crossing bonds are distinct and independent of $\xi$ and $\zeta$.
If $S$ counts successes, its conditional variance is at most $Kc$.
Chebyshev's inequality gives
\begin{equation}
 \mathbb P\left(S<\frac{Kc}{2}\,\middle|\,\zeta\right)
 \le \frac{4}{Kc}\le\frac{4}{K\rho_0}.
 \label{ext:success-count}
\end{equation}

Let $Z$ count successful entrances whose centers satisfy $g_y\le1-\tau$.
For each selected $y$, \eqref{ext:relay-quality} gives
\[
 \mathbb P(O_y,\ g_y\le1-\tau\mid\zeta)
 \le\frac{q^M\gamma}{\tau}.
\]
Indeed, after fixing $\zeta$ the selected vertices are fixed, and
$O_y\cap\{g_y\le1-\tau\}$ is an event of $\xi$, whose law has not
changed. For bonds the additional crossing bit is independent and
contributes a factor $q$. Summing these bounds gives
\[
 \mathbb E[Z\mid\zeta]\le\frac{Kc\gamma}{\tau},
 \qquad
 \mathbb P\left(Z\ge\frac{Kc}{2}\,\middle|\,\zeta\right)
 \le\frac{2\gamma}{\tau}.
\]
If $S\ge Kc/2$ and $Z<Kc/2$, some successful entrance has its center
in $A_\xi$. Thus $o$ reaches $A_\xi$. Integrating over $\zeta$ and
including the low-entry event from \eqref{ext:layer}, we conclude that
\begin{equation}
 \mathbb P(o\leftrightarrow B,\ o\not\leftrightarrow A_\xi)
 \le \frac{1}{\alpha(R-m)}
       +\frac{4}{K\rho_0}+\frac{2\gamma}{\tau}.
 \label{ext:reach-relay}
\end{equation}

\emph{Passing from the relay to the target.}
The estimate \eqref{ext:reach-relay} is now under the original law.
To control the remaining failure event, condition on $\xi$ alone.
Then $A_\xi$ is a deterministic set, the remaining bits have their
product law, and the conditional target probability of each
$v_y\in A_\xi$ is the number $g_y>1-\tau$ used in its definition.
Corollary~\ref{cmp:failure} therefore gives
\[
 \mathbb P(o\leftrightarrow A_\xi,\ o\not\leftrightarrow T
             \mid\xi)\le\tau,
\]
including the trivial case $A_\xi=\varnothing$.
Integrating this bound, and splitting the desired failure event
according to whether $o$ reaches $A_\xi$, yields
\[
 \mathbb P(o\leftrightarrow B,\ o\not\leftrightarrow T)
 \le\frac{1}{\alpha(R-m)}
 +\frac{4}{K\rho_0}+\frac{2\gamma}{\tau}+\tau<\eta.
\]
Each of the four terms is less than $\eta/4$ by
\eqref{prelim:budgets} and \eqref{prelim:scale-inequalities}.
\end{proof}
